% !TeX program = pdflatex
\documentclass[9pt,a4paper]{extarticle}
\usepackage[a4paper,left=10mm,right=10mm,top=9mm,bottom=12mm,footskip=6mm]{geometry}
\usepackage[T1]{fontenc}
\usepackage[utf8]{inputenc}
\usepackage[ngerman]{babel}
\usepackage{lmodern,microtype,amsmath,array,tabularx,xcolor,listings,tcolorbox,fancyhdr,hyperref}
\renewcommand{\familydefault}{\sfdefault}
\definecolor{ink}{HTML}{182C3A}
\definecolor{accent}{HTML}{294F66}
\definecolor{string}{HTML}{345B39}
\definecolor{muted}{HTML}{50545A}
\definecolor{paper}{HTML}{F3F4F5}
\definecolor{rule}{HTML}{A7ADB1}
\color{ink}
\hypersetup{colorlinks=true,urlcolor=accent,pdftitle={Python-Kurzreferenz | Digitalisierung in Unternehmen},pdfauthor={Dr. Christopher Gerling}}
\setlength{\parindent}{0pt}
\setlength{\parskip}{.7pt}
\setlength{\tabcolsep}{3pt}
\renewcommand{\arraystretch}{1.13}
\pagestyle{fancy}\fancyhf{}
\renewcommand{\headrulewidth}{0pt}
\fancyfoot[L]{\fontsize{6.3}{7.5}\selectfont Dr. Christopher Gerling\enspace·\enspace HWR Berlin\enspace·\enspace Digitalisierung in Unternehmen}
\fancyfoot[R]{\fontsize{6.3}{7.5}\selectfont Stand 17.11.2026\enspace |\enspace\thepage\,/\,2}
\lstset{language=Python,basicstyle=\ttfamily\fontsize{7.7}{8.8}\selectfont\color{ink},keywordstyle=\bfseries\color{accent},stringstyle=\color{string},commentstyle=\itshape\color{muted},showstringspaces=false,keepspaces=true,columns=fixed,basewidth=4.04pt,upquote=true,aboveskip=2pt,belowskip=2pt,tabsize=4,breaklines=false,morekeywords={True,False,None},literate={ä}{{\"a}}1 {ö}{{\"o}}1 {ü}{{\"u}}1 {Ä}{{\"A}}1 {Ö}{{\"O}}1 {Ü}{{\"U}}1 {ß}{{\ss}}1}
\newcommand{\py}[1]{\mbox{\ttfamily\fontsize{7.7}{8.8}\selectfont\detokenize{#1}}}
\newtcolorbox{karte}[2][]{before upper=\raggedright,colback=white,colframe=rule,colbacktitle=paper,coltitle=ink,fonttitle=\bfseries\fontsize{9.4}{11}\selectfont,title={#2},boxrule=.45pt,arc=1mm,left=2mm,right=2mm,top=.4mm,bottom=.5mm,toptitle=.8mm,bottomtitle=.8mm,before skip=0pt,after skip=.7mm,#1}
\newcommand{\merk}[1]{\par{\bfseries Merke:} #1\par}
\newenvironment{tabelle}[1]{\tabularx{\linewidth}{@{}#1@{}}}{\endtabularx}
% Beide Kopfzeilen nutzen gemeinsame Grundlinien links und rechts.
\newcommand{\kopf}{%
\noindent\hbox to \linewidth{{\fontsize{19}{21}\selectfont\bfseries Python 3\enspace\mdseries Kurzreferenz}\hfil{\fontsize{8.1}{9.4}\selectfont\bfseries Dr. Christopher Gerling}}\par
\vspace{1mm}
\noindent\hbox to \linewidth{{\fontsize{9.5}{12}\selectfont Digitalisierung in Unternehmen · BWL Bank}\hfil{\fontsize{8.1}{9.4}\selectfont HWR Berlin}}\par
\vspace{2mm}{\color{accent}\hrule height 1pt}\vspace{3mm}}

\begin{document}
\fontsize{8.1}{9.4}\selectfont
\kopf
\begin{minipage}[t][242mm][s]{.489\textwidth}\vspace{0pt}
\begin{karte}{01\enspace Werte, Typen und Wertzuweisung}
\begin{lstlisting}
anzahl = 3                  # int: ganze Zahl
zins = 0.025                # float: Dezimalzahl
name = "Alice"              # str: Zeichenkette
aktiv = True                # bool: Wahrheitswert
werte = [100, 50, 25]       # list: Liste
punkt = (3, 4)              # tuple: Tupel
\end{lstlisting}
\py{type(zins)} ermittelt den Typ: \py{float}.\par
\textbf{Wertzuweisung mit dem Zeichen \py{=}}\par Den Ausdruck rechts auswerten und den Wert dem Variablennamen links zuweisen.
\begin{lstlisting}
tageslimit = 30000
bereits_genutzt = 10000
restlimit = tageslimit - bereits_genutzt
anzahl = anzahl + 1         # bisherigen Wert erhöhen
\end{lstlisting}
Namen: \py{kunde_2}, nicht \py{2kunde}. Groß-/Kleinschreibung zählt. Schlüsselwörter und Namen wie \py{print} nicht als eigene Variablennamen verwenden.
\end{karte}
\vfill
\begin{karte}{02\enspace Eingabe, Konvertierung, Ausgabe}
\begin{lstlisting}
name = input("Name: ")
alter = int(input("Alter: "))
betrag = float(input("Betrag: "))
print("Hallo " + name)
print("Alter: " + str(alter))
\end{lstlisting}
\textbf{\py{input()} liefert immer Text.} Für Rechnungen zunächst in eine Zahl konvertieren. Gültige Zahleneingaben vorausgesetzt; Dezimalpunkt statt Komma.
\begin{tabelle}{l X}
\py{int("12")} & \py{12}: ganze Zahl \\
\py{int(12.9)} & \py{12}: Nachkommastellen abschneiden \\
\py{float("12.5")} & \py{12.5}: Fließkommazahl \\
\py{str(12)} & \py{"12"}: Zeichenkette \\
\end{tabelle}
\begin{lstlisting}
print("12" + "3")           # 123: Text verbinden
print(12 + 3)               # 15: Zahlen addieren
print("Betrag: " + str(12.5))
\end{lstlisting}
Bei \py{+} entweder \textbf{zwei Strings verketten} oder \textbf{zwei Zahlen addieren}.\par Zahlen für Textausgaben mit \py{str(...)} konvertieren.\par Leerzeichen in den String schreiben, z.\,B. \py{"Hallo "}.
\end{karte}
\vfill
\begin{karte}{03\enspace Rechnen mit Zahlen}
\begin{tabelle}{l X l}
\textbf{Operator} & \textbf{Bedeutung} & \textbf{Beispiel} \\
\py{+ - *} & plus, minus, mal & \py{7 * 3} $\to$ \py{21} \\
\py{/} & Division & \py{6 / 3} $\to$ \py{2.0} \\
\py{//} & Quotient, abgerundet & \py{7 // 3} $\to$ \py{2} \\
\texttt{\%} & Rest (Modulo) & \texttt{7 \% 3} $\to$ \py{1} \\
\py{**} & Potenz & \py{2 ** 3} $\to$ \py{8} \\
\end{tabelle}
\textbf{Klammern zuerst}, dann Potenzen, Punkt- und Strichrechnung. Zahlen können auch in Variablen stehen; mit ihnen wird genauso gerechnet.
\end{karte}
\vfill
\begin{karte}{04\enspace Strings bearbeiten}
\py{"..."} und \py{'...'} kennzeichnen Zeichenketten.
\begin{lstlisting}
text = "HWR Berlin"
text.upper()                # "HWR BERLIN"
text.lower()                # "hwr berlin"
text.replace("HWR", "Uni")  # "Uni Berlin"
text.split()                # ["HWR", "Berlin"]
\end{lstlisting}
\textbf{Strings bleiben unverändert.} Diese Methoden liefern einen neuen String (\py{split()} eine Liste) als \textbf{Rückgabewert}. Sie geben selbst nichts aus. Anzeigen mit \py{print(...)} oder als letztes Ergebnis einer Notebookzelle.\par
\textbf{Zur dauerhaften Speicherung} das Ergebnis wieder der Variablen zuweisen:
\begin{lstlisting}
text = text.lower()
text                        # "hwr berlin"
\end{lstlisting}
\end{karte}
\end{minipage}\hfill
\begin{minipage}[t][242mm][s]{.489\textwidth}\vspace{0pt}
\begin{karte}{05\enspace Vergleiche und Aussagenlogik}
\begin{tabelle}{l X}
\py{==} / \py{!=} & gleich / ungleich \\
\py{< <= > >=} & kleiner, kleiner gleich, größer, größer gleich \\
\py{x in folge} & enthalten? Auch Teiltext: \py{"B" in "Bank"} \\
\py{x not in folge} & nicht enthalten? \\
\end{tabelle}
\smallskip
\textbf{Nicht verwechseln:} \py{=} weist einen Wert zu; \py{==} prüft, ob zwei Werte gleich sind.\par\smallskip
\begin{tabularx}{\linewidth}{@{}XX|XXX@{}}
\textbf{a} & \textbf{b} & \textbf{not a} & \textbf{a and b} & \textbf{a or b} \\
False & False & True & False & False \\
False & True & True & False & True \\
True & False & False & False & True \\
True & True & False & True & True \\
\end{tabularx}
\smallskip
\py{not} kehrt um; \py{and} verlangt zwei wahre Aussagen.\par
\textbf{\py{or} ist ein inklusives Oder:} Mindestens eine Aussage ist wahr; beide dürfen wahr sein.\par
\textbf{Reihenfolge:} zuerst Klammern, dann \py{not}, \py{and}, \py{or}. Vergleiche bilden dabei die Teilaussagen.\par\smallskip
Beide Seiten vollständig formulieren:\par
\textbf{Richtig}\quad \py{x == 5 or x == 7}\par
\textbf{Falsch}\quad \py{x == 5 or 7}
\end{karte}
\vfill
\begin{karte}{06\enspace Verzweigungen: if und else}
Eine \textbf{Bedingung} ist z.\,B. ein Vergleich oder eine Verknüpfung aus Kasten 05. Hier sei \py{saldo = -20}.\par\smallskip
\begin{lstlisting}
if saldo < 0:
    print("negativ")
else:
    print("ab null")
\end{lstlisting}
\textbf{Doppelpunkt + Einrückung} bilden einen Block. Üblich sind vier Leerzeichen pro Ebene. \py{else} steht auf derselben Ebene wie das zugehörige \py{if}.\par
Bedingung wahr: \py{if}-Block; sonst: \py{else}-Block. Ohne \py{else} geht es bei falscher Bedingung nach dem Block weiter. Einzelne \py{if} werden unabhängig geprüft.
\end{karte}
\vfill
\begin{karte}{07\enspace Zeichenketten prüfen}
\begin{tabelle}{l X}
\py{text.isalpha()} & nur Buchstaben, auch Umlaute \\
\py{text.isalnum()} & nur Buchstaben / numerische Zeichen \\
\py{text.isdigit()} & nur Ziffern \\
\py{text.islower()} & alle Buchstaben klein \\
\py{text.isupper()} & alle Buchstaben groß \\
\end{tabelle}
Die ersten drei verlangen mindestens ein Zeichen; \py{islower()} und \py{isupper()} mindestens einen Buchstaben mit Groß-/Kleinschreibung. Leere Strings liefern \py{False}.\par
\textbf{Leerzeichen zählen hier als Sonderzeichen:} Sie sind weder Buchstaben noch Ziffern; \py{isalnum()} liefert dann \py{False}. Dasselbe gilt für Satzzeichen.
\begin{lstlisting}
text = "HWR Berlin"
print(text.isalnum())       # False: Leerzeichen
print("HWR123".isalnum())   # True
print("Berlin".isalpha())   # True
print("HWR 123".isupper())  # True
\end{lstlisting}
\py{isupper()} ignoriert z.\,B. Leerzeichen und Ziffern; es prüft die Groß-/Kleinschreibung der Buchstaben.
\end{karte}
\vfill
\begin{karte}{08\enspace Kommentare}
\texttt{\#} beginnt außerhalb eines Strings einen Kommentar. Python ignoriert den Rest dieser Zeile.
\begin{lstlisting}
betrag = 120                # Rechnungsbetrag in Euro
\end{lstlisting}
Kommentare erklären den Code oder zeigen hier das erwartete Ergebnis. Sie werden nicht ausgegeben.
\end{karte}
\end{minipage}
\newpage
\renewcommand{\arraystretch}{1.0}
\fontsize{7.7}{8.8}\selectfont
\lstset{basicstyle=\ttfamily\fontsize{7.4}{8.2}\selectfont\color{ink}}
\kopf
\begin{minipage}[t][255mm][s]{.489\textwidth}\vspace{0pt}
\begin{karte}{09\enspace Folgen: Indexierung und Slicing}
\textbf{Strings, Listen und Tupel sind geordnete Folgen.} Ihre Elemente sind über Positionen (Indizes) erreichbar.
\begin{center}
\begin{tikzpicture}[x=8mm,y=5mm,font=\sffamily\fontsize{8}{9}\selectfont]
\node[anchor=east] at (-.65,0) {\py{text =}};
\foreach \i/\ch/\neg in {0/B/-6,1/e/-5,2/r/-4,3/l/-3,4/i/-2,5/n/-1}{
\node[draw=rule,fill=paper,minimum width=7.6mm,minimum height=6mm,font=\ttfamily\bfseries] at (\i,0) {\ch};
\node at (\i,-1) {\i};\node at (\i,-1.8) {\neg};}
\node[anchor=east] at (-.65,-1) {Index};
\node[anchor=east] at (-.65,-1.8) {Negativer Index};
\end{tikzpicture}
\end{center}\vspace{-2mm}
\begin{lstlisting}
text = "Berlin"
text[0]                     # "B"
text[-1]                    # "n"

werte = [10, 20, 30, 40, 50]
werte[1]                    # 20
\end{lstlisting}
\textbf{Ab 0 zählen.} \py{-1} bezeichnet das letzte Element; sein positiver Index ist \py{len(folge)-1}.\par\smallskip
\textbf{Slicing mit \py{folge[start:ende]}}\par Start enthalten, Ende ausgeschlossen.\par Die Grenzen liegen zwischen den Elementen:
\begin{center}
\begin{tikzpicture}[x=8mm,y=5mm,font=\sffamily\fontsize{8}{9}\selectfont]
\node[anchor=east] at (-.7,0) {\py{text =}};
\foreach \i/\ch in {0/B,1/e,2/r,3/l,4/i,5/n}{
\ifnum\i>0\ifnum\i<4
\fill[accent!20] (\i-.5,-.5) rectangle (\i+.5,.5);
\fi\fi
\node[font=\ttfamily\bfseries] at (\i,0) {\ch};}
\foreach \j in {0,...,6}{\draw[rule] (\j-.5,-.5)--(\j-.5,.5);\node at (\j-.5,-1) {\j};}
\draw[ink,thick] (.5,-.5)--(.5,.5);\draw[ink,thick] (3.5,-.5)--(3.5,.5);
\node[anchor=east] at (-.7,-1) {Grenze};
\end{tikzpicture}
\end{center}\vspace{-2mm}
\begin{lstlisting}
text[1:4]                   # "erl"
text[:3]                    # "Ber"
text[3:]                    # "lin"
text[::2]                   # "Bri"
text[::-1]                  # "nilreB"
werte[1:3]                  # [20, 30]
\end{lstlisting}
\py{folge[start:ende:schritt]} ergänzt die Schrittweite. Ohne Grenzen: gesamte Folge; ohne Schritt: 1.
\end{karte}
\vfill
\begin{karte}{10\enspace Generische Operationen auf Folgen}
\py{folge} steht für einen String, eine Liste oder ein Tupel. Die Elemente müssen zur Operation passen; zum Vergleichen und Sortieren müssen sie vergleichbar sein.
\begin{tabelle}{l X}
\py{len(folge)} & Anzahl der Elemente / Zeichen \\
\py{x in folge} & Mitgliedschaft; bei Strings auch Teiltext \\
\py{folge.count(x)} & Anzahl der Vorkommen \\
\py{folge.index(x)} & erste Position; Wert muss vorkommen \\
\py{min(folge)} & kleinster Wert; Folge nicht leer \\
\py{max(folge)} & größter Wert; Folge nicht leer \\
\py{sorted(folge)} & neue sortierte \textbf{Liste} \\
\py{sum(folge)} & Summe \textbf{numerischer} Elemente \\
\end{tabelle}

\end{karte}
\vfill
\begin{karte}{11\enspace Listen anlegen und verändern}
\begin{lstlisting}
zahlen = []                 # leere Liste
zahlen = [7, 2, 5]          # Liste mit drei Elementen
\end{lstlisting}
\textbf{Verändern}\par Die Liste wird direkt geändert; die Methoden liefern \py{None}.\par\smallskip
\begingroup\renewcommand{\arraystretch}{.95}
\begin{tabelle}{l X}
\py{zahlen[i] = x} & Element an Index \py{i} ersetzen \\
\py{zahlen.append(x)} & ein Element hinten anhängen \\
\py{zahlen.insert(i, x)} & vor Index \py{i} einfügen \\
\py{zahlen.remove(x)} & ersten passenden Wert entfernen \\
\py{del zahlen[i]} & Element an Index \py{i} löschen \\
\py{zahlen.sort()} & Liste selbst aufsteigend sortieren \\
\py{zahlen.reverse()} & Reihenfolge der Liste umkehren \\
\end{tabelle}\par\smallskip
\textbf{Abfragen}\par Die Liste bleibt unverändert; das Ergebnis ist eine Zahl.\par
\begin{tabelle}{l X}
\py{zahlen.index(x)} & erste Position von \py{x} (muss vorkommen) \\
\py{zahlen.count(x)} & Anzahl der Vorkommen von \py{x} \\
\end{tabelle}\endgroup
\par\smallskip
Ergebnisse mit \py{print(...)} anzeigen.\par
\py{sorted(zahlen)} liefert eine neue Liste; das Original bleibt erhalten.
\end{karte}

\vfill
\begin{karte}{12\enspace Tabellen: Liste von Zeilen}
Spalten: Sender, Empfänger, Betrag, Währung.
\begin{lstlisting}
tx = [["Alice", "Bob", 100.0, "EUR"],
      ["Fred", "Bob", 18.36, "EUR"]]
tx[0][2]                    # 100.0
tx[0][2] = 125.0            # Zelle ändern
tx.append(["Harry", "Bob", 50.0, "USD"])
del tx[-1]                  # letzte Zeile entfernen
\end{lstlisting}
\py{tabelle[zeile][spalte]} wählt erst die Zeile, dann die Spalte; beide ab 0. Eine Kopfzeile wird mitgezählt.
\end{karte}
\end{minipage}\hfill
\begin{minipage}[t][255mm][s]{.489\textwidth}\vspace{0pt}
\begin{karte}{13\enspace for-Schleife: über Elemente laufen}
\begin{lstlisting}
banken = ["ING", "CB", "DB"]
for kuerzel in banken:
    print(kuerzel)
\end{lstlisting}
\py{banken} ist die Liste; \py{kuerzel} steht nacheinander für \py{"ING"}, \py{"CB"} und \py{"DB"}. Der Block läuft \textbf{einmal pro Element}.\par Durchlaufbare Objekte heißen \textbf{Iterables}: z.\,B. Listen, Tupel, Strings und \py{range}.\par
\textbf{Zahlen durchlaufen mit \py{range}}
\begin{tabelle}{l X}
\py{range(5)} & 0, 1, 2, 3, 4 \\
\py{range(1, 6)} & 1, 2, 3, 4, 5 \\
\py{range(0, 11, 2)} & 0, 2, 4, 6, 8, 10 \\
\py{range(5, 0, -1)} & 5, 4, 3, 2, 1 \\
\end{tabelle}
\begin{lstlisting}
for i in range(1, 4):
    print(i)                # 1, 2, 3 untereinander
\end{lstlisting}
\py{range(start, ende, schritt)} schließt das Ende aus.\par Die Schrittweite darf nicht 0 sein.\par Für 1 bis einschließlich \py{n} gilt \py{range(1, n + 1)}.
\end{karte}
\vfill
\begin{karte}{14\enspace while-Schleife: bedingt wiederholen}
\begin{lstlisting}
i = 1                       # Startwert
while i <= 5:               # vor jedem Durchlauf prüfen
    print(i)                # 1 bis 5 untereinander
    i = i + 1               # Zähler fortschreiben
print("Ende")
\end{lstlisting}
\textbf{Solange die Bedingung wahr ist}, wird wiederholt.\par Die Bedingung wird vor jedem Durchlauf geprüft.\par Im Beispiel läuft der Block jeweils vollständig durch.\par Den Zähler vorab setzen und im Block fortschreiben.
\end{karte}
\vfill
\begin{karte}{15\enspace Abbrechen und Überspringen}
\begin{minipage}[t]{.48\linewidth}\vspace{0pt}
\textbf{Schleife beenden}
\begin{lstlisting}
for c in "bank":
    if c == "n":
        break
    print(c)
\end{lstlisting}
Ausgabe: b, a
\end{minipage}\hfill
\begin{minipage}[t]{.48\linewidth}\vspace{0pt}
\textbf{Durchlauf verkürzen}
\begin{lstlisting}
for c in "bank":
    if c == "n":
        continue
    print(c)
\end{lstlisting}
Ausgabe: b, a, k
\end{minipage}\par\smallskip
Jeweils ein Buchstabe pro Ausgabezeile.\par
\textbf{\py{break}} beendet die innerste Schleife sofort.\par
\textbf{\py{continue}} überspringt den Rest dieses Durchlaufs.\par
Bei \py{while} den Zähler auch bei \py{continue} fortschreiben, sonst droht eine Endlosschleife.
\end{karte}
\vfill
\begin{karte}{16\enspace Einrückung und Verschachtelung}
\textbf{Verschachtelung:} ein Block innerhalb eines anderen. Pro Ebene \textbf{vier Leerzeichen} mehr einrücken.\par\smallskip
\begin{tikzpicture}[x=4.04pt,y=11pt]
\draw[rule,line width=.6pt] (1.4,-.35)--(1.4,-3.2);
\draw[accent,line width=.6pt] (5.4,-1.35)--(5.4,-2.2);
\node[anchor=base west,inner sep=0pt] at (0,0) {\py{for zahl in [1, 2]:}};
\node[anchor=base west,inner sep=0pt] at (4,-1) {\py{if zahl > 1:}};
\node[anchor=base west,inner sep=0pt] at (8,-2) {\py{print(zahl)}};
\node[anchor=base west,inner sep=0pt] at (4,-3) {\py{print("weiter")}};
\node[anchor=base west,inner sep=0pt] at (0,-4) {\py{print("Ende")}};
\node[anchor=base west,inner sep=0pt,font=\scriptsize\color{muted}] at (32,-1) {4 Leerzeichen};
\node[anchor=base west,inner sep=0pt,font=\scriptsize\color{muted}] at (32,-2) {8 Leerzeichen};
\end{tikzpicture}\par\smallskip
\py{print("weiter")} liegt außerhalb des \py{if}-Blocks, aber in der Schleife; \py{print("Ende")} nach der Schleife.\par Die Hilfslinien gehören nicht zum Code.
\end{karte}
% Beide letzten Themen stehen untereinander am Ende der rechten Spalte.
\vfill
\begin{karte}{17\enspace Eigene Funktionen}
\begin{lstlisting}
def endkapital(k0, zins, jahre):
    return k0 * (1 + zins) ** jahre
betrag = endkapital(1000, 0.07, 10)
\end{lstlisting}
\py{def} definiert; erst der \textbf{Aufruf} führt aus.\par \py{k0}, \py{zins} und \py{jahre} sind \textbf{Parameter}.\par Beim Aufruf erhalten sie der Reihe nach die Werte \py{1000}, \py{0.07} und \py{10}.\par \py{return} beendet die Funktion und liefert das Ergebnis; hier wird es in \py{betrag} gespeichert. \py{print()} zeigt nur an.
\end{karte}
\vfill
\begin{karte}{18\enspace Bibliotheken verwenden}
\begin{lstlisting}
import math
math.sqrt(256)              # 16.0
math.ceil(2.3)              # 3
\end{lstlisting}
\py{import} lädt ein Modul. Zugriff über \py{modul.name}.\par \py{sqrt} berechnet die Wurzel; \py{ceil} rundet ganzzahlig auf.
\end{karte}

\end{minipage}
\end{document}
